\chapter{机器学习}

\gls{ML}旨在使模型从观测数据中学习潜在的规律，并利用所学习的规律完成预测或结构发现等任务。设一项研究共观测到$n$个样本，每个样本由$p$个\gls{Feature}描述。记第$i$个样本的特征向量为：
\begin{equation*}
	x_i=(x_{i1},x_{i2},\ldots,x_{ip})^\top\in X
\end{equation*}
其中$x_{ij}$表示第$i$个样本在第$j$个特征上的取值，称$X$为特征空间。\par
根据观测数据是否包含需要预测的目标信息，机器学习方法可以大致分为\gls{SupervisedLearning}和\gls{UnsupervisedLearning}。\par
在监督学习中，每个特征向量$x_i$都对应一个输出结果$y_i\in Y$，其中$y_i$称为第$i$个样本的响应变量，$Y$称为输出空间。监督学习的观测数据集记为：
\begin{equation*}
	\mathcal{D}=\{(x_i,y_i)\}_{i=1}^n
\end{equation*}
将所有样本的输入向量按行排列，可以得到特征矩阵：
\begin{equation*}
	X=
	\begin{pmatrix}
		x_1^\top \\
		x_2^\top \\
		\vdots \\
		x_n^\top
	\end{pmatrix}
	=
	\begin{pmatrix}
		x_{11}&x_{12}&\cdots&x_{1p} \\
		x_{21}&x_{22}&\cdots&x_{2p} \\
		\vdots&\vdots&\ddots&\vdots \\
		x_{n1}&x_{n2}&\cdots&x_{np}
	\end{pmatrix}
\end{equation*}
相应的输出向量记为：
\begin{equation*}
	y=(y_1,y_2,\dots,y_n)^\top
\end{equation*}
监督学习的目标是利用观测数据集$\mathcal{D}$学习一个从特征空间$X$到输出空间$Y$的预测函数：
\begin{equation*}
	f:X\longrightarrow Y
\end{equation*}
使得对于观测数据中未出现的新输入$x\in X$，模型能够给出相应的预测$f(x)$。\par
例如，考虑利用血常规检查结果辅助判断受检者是否患有某种疾病。对于每位既往患者，我们既记录其白细胞计数、中性粒细胞比例、血红蛋白浓度和血小板计数等血常规指标，又记录医生结合临床表现、血常规结果及其他检查所确定的诊断结果。于是，每位患者的血常规指标构成一个特征向量$x_i$，其临床诊断结果构成相应的输出变量$y_i$。收集一批既往患者的数据后，可以得到监督学习的观测数据集$\mathcal{D}=\{(x_i,y_i)\}_{i=1}^n$。模型利用这些已知诊断结果的历史数据，学习血常规指标与患病状态之间的潜在关系。当获得一位新受检者的血常规数据时，便可以将其作为新的输入，由模型预测该受检者是否可能患病，从而为医生的临床判断提供辅助信息。表\ref{tab:supervised-learning-example}给出了一个简化的观测数据集示例。
\begin{table}[H]
	\centering
	\caption{基于血常规指标判断是否患病的监督学习数据示例}
	\label{tab:supervised-learning-example}
	\begin{tabular}{cccccc}
		\toprule
		样本
		&白细胞计数
		&中性粒细胞比例
		&血红蛋白
		&血小板计数
		&是否患病 \\
		\midrule
		$1$ & $6.2$ & $58$ & $145$ & $245$ & 否 \\
		$2$ & $12.8$ & $82$ & $132$ & $310$ & 是 \\
		$3$ & $5.7$ & $54$ & $151$ & $228$ & 否 \\
		$4$ & $14.1$ & $86$ & $126$ & $335$ & 是 \\
		$5$ & $7.0$ & $61$ & $139$ & $260$ & 否 \\
		$6$ & $11.6$ & $79$ & $128$ & $298$ & 是 \\
		\bottomrule
	\end{tabular}
\end{table}\par
在无监督学习中，观测数据不包含预先给定的输出变量，因而观测数据集记为：
\begin{equation*}
	\mathcal{D}=\{x_i\}_{i=1}^n
\end{equation*}
无监督学习的目标不是预测某个已知的输出，而是从输入数据中发现潜在的结构或模式。例如，可以根据样本之间的相似性将其划分为若干组，或者使用较少的变量概括原始数据中的主要信息。\par
在监督学习中，根据输出空间$Y$的类型，可以进一步将问题分为分类问题和回归问题。\par
当响应变量只取有限个离散类别时，该问题称为\textbf{分类问题}。设所有可能的类别为$Y=\{1,2,\dots,K\}$，其中$K$为类别数。分类模型需要根据特征判断一个样本所属的类别。例如，根据邮件内容判断其是否为垃圾邮件，或者根据医学检查结果判断患者是否患病。此时我们又称第$i$个样本的响应变量$y_i$为其\gls{Label}。\par
当响应变量为连续数值时，该问题称为\textbf{回归问题}，通常有$Y\subseteq\mathbb{R}$。回归模型需要学习特征与数值型输出之间的关系。例如，根据房屋面积和位置预测房价，或者根据患者的临床指标预测某项连续的生理数值。\par
分类问题在形式上也可以转化为回归问题：模型首先根据样本特征输出一个连续的分类得分，再根据预先设定的阈值确定样本所属的类别。例如，可以将分类得分大于$0$的样本判为一类，将分类得分小于或等于$0$的样本判为另一类。

\section{模型评价}

在监督学习中，模型评价是衡量模型预测能力、比较不同模型并指导模型选择的重要环节。

\subsection{分类任务}
分类问题可以根据类别数划分为二分类问题与多分类问题，二分类问题即$K=2$的分类问题，多分类问题对应$K\geqslant3$的分类问题。
\subsubsection{二分类}
在二分类问题中，通常将研究中重点关注、希望模型识别出的类别称为\textbf{正类}，将另一个类别称为\textbf{负类}。\par
\begin{definition}
	\gls{ConfusionMatrix}按真实标签与预测标签的组合统计样本个数，其每个元素含义如下：
	\begin{itemize}
		\item \gls{TP}：真实为正类且预测为正类的样本数；
		\item \gls{TN}：真实为负类且预测为负类的样本数；
		\item \gls{FP}：真实为负类但预测为正类的样本数；
		\item \gls{FN}：真实为正类但预测为负类的样本数。
	\end{itemize}
\end{definition}
\begin{table}[H]
	\centering
	\caption{二分类问题的混淆矩阵}
	\begin{tabular}{ccc}
		\toprule
		真实类别$\backslash$预测类别
		&正类($\hat{y}=1$)
		&负类($\hat{y}=0$) \\
		\midrule
		正类($y=1$)
		&TP
		&FN \\
		负类($y=0$)
		&FP
		&TN \\
		\bottomrule
	\end{tabular}
\end{table}
\begin{definition}
  \gls{Accuracy}度量模型预测正确的比例：
  \begin{equation*}
    \operatorname{Accuracy}=\frac{\operatorname{TP}+\operatorname{TN}}{\operatorname{TP}+\operatorname{TN}+\operatorname{FP}+\operatorname{FN}}
  \end{equation*}
\end{definition}
Accuracy 将正类和负类的正确预测汇总为一个整体比例，因而默认两类样本具有相同的重要性，并且不同类型的分类错误具有相同的代价。然而在许多实际问题中，假正类和假负类造成的损失并不相同，我们也可能更关注模型在某一特定类别上的预测表现。\par
例如，在贷款违约预测中，若将违约客户定义为正类，则将实际违约的客户预测为不违约会使银行承担贷款损失，而将实际不违约的客户预测为违约主要意味着失去一次潜在的放贷机会。这两类错误的经济代价并不相同。因此，即使两个模型具有相同的 Accuracy，它们识别违约客户的能力也可能存在显著差异。为了考察模型针对特定类别的表现，需要进一步引入别的指标。
\begin{definition}
  \gls{Precision}度量被预测为正类的样本中有多少是真正的正类：
  \begin{equation*}
    \operatorname{Precision}=\frac{\operatorname{TP}}{\operatorname{TP}+\operatorname{FP}}
  \end{equation*}
\end{definition}
\begin{note}
  精确率关心正类预测结果的准确性，在系统对错误正类具有高成本的任务（如垃圾邮件过滤）中尤为重要。较高的精确率意味着模型产生的误报较少。
\end{note}
\begin{definition}
  \gls{Recall}又称\gls{Sensitivity}，衡量真实正类中有多少被模型成功识别，定义为：
  \begin{equation*}
    \operatorname{Recall}=\frac{\operatorname{TP}}{\operatorname{TP}+\operatorname{FN}}
  \end{equation*}
\end{definition}
\begin{note}
  召回率关注对正类的漏判率，适用于需要尽可能检出正类的场景（如疾病筛查或欺诈检测）。高召回率意味着漏报较少，但可能伴随较多假阳性，需要结合精确率来进行分析。
\end{note}
\begin{definition}
  \gls{F1Score}是精确率与召回率的调和平均，用于在二者之间取得平衡：
  \begin{equation*}
    \operatorname{F1}=2\times\frac{\operatorname{Precision}\times\operatorname{Recall}}{\operatorname{Precision}+\operatorname{Recall}}
  \end{equation*}
\end{definition}
\begin{note}
  调和平均的性质使$\operatorname{F1}$得分对低值敏感——当精确率或召回率其中之一很低时，$\operatorname{F1}$会显著下降，因此它常用作综合评价指标，适用于精确率和召回率同等重要的情形。
\end{note}
\begin{definition}
	$\operatorname{F}_{\beta}$分数是$\operatorname{F1}$得分的推广，通过参数$\beta>0$调节精确率与召回率的重要性：
	\begin{equation*}
		\operatorname{F}_{\beta}=(1+\beta^2)\frac{\operatorname{Precision}\times\operatorname{Recall}}{\beta^2\operatorname{Precision}+\operatorname{Recall}}=\dfrac{1}{\dfrac{\beta^2}{1+\beta^2}\dfrac{1}{\operatorname{Recall}}+\dfrac{1}{1+\beta^2}\dfrac{1}{\operatorname{Precision}}}
	\end{equation*}
\end{definition}
\begin{note}
	当$|\beta|=1$时，$\operatorname{F}_{\beta}$退化为$\operatorname{F1}$；当$|\beta|>1$时，召回率具有更大的权重，适用于更关注漏检的任务；当$0<|\beta|<1$时，精确率具有更大的权重，适用于更关注误报的任务。因此，相比固定地给予精确率与召回率相同权重的$\operatorname{F1}$，$\operatorname{F}_{\beta}$能够根据具体任务中两类错误的重要程度进行调整。
\end{note}
许多分类器并不直接输出类别，而是为每个样本输出一个连续得分或属于正类的概率，并在给定阈值后将其转化为分类结果。然而，分类阈值并不存在适用于所有任务的固定选择，即使模型输出的是概率，也不能机械地将$0.5$作为阈值。阈值的选择取决于正类的基础发生率、假正类与假负类的相对代价以及模型所服务的具体决策目标。\par
以罕见病筛查为例，假设某种疾病在人群中的患病率仅为$0.1\%$。若模型预测某位受检者的患病概率为$10\%$，那么其患病风险已经达到普通人群的$100$ 倍，显然应当被识别为高风险个体并接受进一步检查。然而，$10\%$仍然低于$50\%$；若机械地将$0.5$作为分类阈值，该受检者就会被判定为阴性。类似地，在身份认证或银行人脸识别中，系统可能更重视防止冒用身份，因此会采用更加严格的判定阈值，以降低错误接受非本人的概率。\par
由此可见，同一个模型在不同应用场景中可能需要采用完全不同的阈值，而模型在某一个固定阈值下的表现不能完整反映其分类能力。为了减少评价结果对特定阈值选择的依赖，可以考察阈值连续变化时模型性能的变化，并将其表示为性能曲线。我们希望一个模型在较广的阈值范围内均能保持较好的正负类区分能力，因此可以进一步使用曲线下面积评价模型的综合表现。
\begin{definition}
	称：
	 \begin{equation*}
		\operatorname{TPR}=\frac{\operatorname{TP}}{\operatorname{TP}+\operatorname{FN}},\quad
		\operatorname{FPR}=\frac{\operatorname{FP}}{\operatorname{FP}+\operatorname{TN}}
	\end{equation*}
	分别为\gls{TPR}与\gls{FPR}。
\end{definition}
\begin{definition}
  设分类器将得分不低于阈值$t$的样本预测为正类。\gls{ROC}为在阈值取值范围内连续变化$t$，对每个$t$计算对应的：
	\begin{equation*}
		\operatorname{TPR}(t)=\frac{\operatorname{TP}(t)}{\operatorname{TP}(t)+\operatorname{FN}(t)},\quad
		\operatorname{FPR}(t)=	\frac{\operatorname{FP}(t)}{\operatorname{FP}(t)+\operatorname{TN}(t)}
	\end{equation*}
	然后以$\mathrm{FPR}$为横轴、$\mathrm{TPR}$为纵轴作图得到的曲线。\gls{AUROC}定义为ROC曲线与$\mathrm{FPR}$轴和$\mathrm{FPR}=1$围成区域的面积。
\end{definition}
\begin{note}
  ROC曲线反映了阈值变化对$\operatorname{TPR}$与$\operatorname{FPR}$的权衡。$\operatorname{AUROC}$独立于具体阈值，是比较不同模型总体区分能力的常用指标，取值范围通常在$[0.5,1]$，$0.5$表示随机猜测（对任意的$t\in[0,1]$有$\operatorname{FPR}(t)=\operatorname{TPR}(t)$，模型判断正类为正类的概率和模型判断负类为正类的概率一致），$1$表示完美分类。
  \begin{figure}[H]
  	\centering
  	\begin{tikzpicture}[scale=3]
  		% 坐标轴
  		\draw[->] (0,0) -- (1.05,0) node[right] {\scriptsize FPR};
  		\draw[->] (0,0) -- (0,1.05) node[above] {\scriptsize TPR};
  		% 随机分类器的对角线
  		\draw[color=gray, dashed] (0,0) -- (1,1);
  		% 绘制示意 ROC 曲线
  		\draw[color=blue, thick] (0,0) .. controls (0.15,0.5) and (0.3,0.75) .. (1,1);
  		% 填充面积
  		\path[blue!20] (0,0) -- (0.15,0.5) .. controls (0.15,0.5) and (0.3,0.75) .. (1,1) -- (1,0) -- cycle;
  	\end{tikzpicture}
  	\caption{ROC曲线示意}
  \end{figure}
\end{note}
\begin{definition}
 	设模型将得分不低于阈值 $t$ 的样本预测为正类。\gls{PRC}为在阈值取值范围内内连续变化$t$，对每个$t$计算对应的$\operatorname{Precision}(t)$和$\operatorname{Recall}(t)$，然后以$\mathrm{Recall}$为横轴、$\mathrm{Precision}$为纵轴作图得到的曲线。\gls{AUPRC}定义为PRC曲线与$\operatorname{Recall}$轴之间区域的面积。
\end{definition}
\begin{note}
  相比$\operatorname{AUROC}$，$\operatorname{AUPRC}$更关注正类的表现。较高的$\operatorname{AUPRC}$表明模型性能较好。
  \begin{figure}[H]
  	\centering
  	\begin{tikzpicture}[scale=3]
  		% 坐标轴
  		\draw[->] (0,0) -- (1.05,0) node[right] {\scriptsize Recall};
  		\draw[->] (0,0) -- (0,1.05) node[above] {\scriptsize Precision};
  		% 绘制示意 PR 曲线
  		\draw[color=blue, thick] (0,1) .. controls (0.3,0.9) and (0.6,0.6) .. (1,0.2);
  		% 填充面积
  		\path[blue!20] (0,1) .. controls (0.3,0.9) and (0.6,0.6) .. (1,0.2) -- (1,0) -- (0,0) -- cycle;
  	\end{tikzpicture}
  	\caption{PR曲线示意}
  \end{figure}
\end{note}
对于能够输出预测概率的分类模型，模型性能至少涉及两个不同方面：一方面，模型应具有良好的排序能力，即正类样本通常获得高于负类样本的预测分数；另一方面，模型给出的预测概率还应具有可靠的概率解释。前者通常可以通过AUROC评价，但良好的排序能力并不意味着预测概率本身准确。例如，某模型可能始终给真正的正类样本赋予高于负类样本的预测概率，从而获得较高的AUROC，但对于所有预测概率约为$0.8$的样本，实际只有约$60\%$属于正类。此时模型虽然能够较好地区分样本的相对风险高低，却系统性高估了事件发生的绝对概率。因此，还需要进一步考察预测概率与实际事件发生频率之间的一致程度。
\begin{definition}
	设分类模型对样本$x_i$给出的正类预测概率为$\hat{p}_i$。将预测概率的取值范围划分为若干区间$\seq{B}{K}$，对于每个区间$B_k$，分别计算其中样本被预测为正类的概率的平均值以及真实正类比例：
	\begin{equation*}
		\operatorname{Pred}(B_k)=\frac{1}{|B_k|}\sum_{i:\hat{p}_i\in B_k}\hat{p}_i,\quad\operatorname{Obs}(B_k)=\frac{1}{|B_k|}\sum_{i:\hat{p}_i\in B_k}y_i
	\end{equation*}
	其中$|B_k|$表示落入第$k$个概率区间$B_k$中的样本数量。\gls{CalibrationCurve}为以$\operatorname{Pred}(B_k)$为横轴、$\operatorname{Obs}(B_k)$为纵轴作图得到的曲线。
\end{definition}
\begin{note}
	理想情况下校准曲线应与直线$y=x$重合。若校准曲线位于$y=x$下方，则说明模型给出的预测概率整体偏高，反之模型给出的预测概率整体偏低。
	\begin{figure}[H]
		\centering
		\begin{tikzpicture}[scale=3]
			% 坐标轴
			\draw[->] (0,0) -- (1.05,0) node[right] {\scriptsize Predicted Probability};
			\draw[->] (0,0) -- (0,1.05) node[above] {\scriptsize Observed Frequency};
			
			% 完美校准线
			\draw[dashed, gray] (0,0) -- (1,1);
			
			% 示意校准曲线
			\draw[color=blue, thick]
			(0.05,0.08)
			.. controls (0.25,0.18) and (0.40,0.32) ..
			(0.55,0.45)
			.. controls (0.70,0.55) and (0.85,0.70) ..
			(0.98,0.82);
		\end{tikzpicture}
		\caption{校准曲线示意}
	\end{figure}
\end{note}
\subsubsection{多分类}
与二分类问题不同，多分类问题中不存在唯一的正类与负类，因此通常需要先分别考察模型在每个类别上的表现，再通过适当的平均方式得到整体评价指标。\par
对于任意类别$k\in\{1,2,\dots,K\}(K\geqslant3)$，可将类别$k$视为正类，将其余类别合并为负类，并按照二分类问题中的定义计算$\operatorname{TP}_k$、$\operatorname{FP}_k$、$\operatorname{FN}_k$和$\operatorname{TN}_k$，进而得到该类别的$\operatorname{Precision}_k$、$\operatorname{Recall}_k$和$\operatorname{F1}_k$。这种处理方式称为one-versus-rest。\par
分别报告每个类别的评价指标能够完整反映模型在各类别上的表现，但当类别数量较多或需要使用单一数值比较模型时，需要进一步对各类别的指标进行汇总。常用的汇总方式包括Macro平均、Weighted平均和Micro平均。\par
\begin{definition}
	设$m_k$表示类别$k$的某个评价指标，该指标的Macro平均定义为：
	\begin{equation*}
		m_{\mathrm{macro}}=\frac{1}{K}\sum_{k=1}^{K}m_k
	\end{equation*}
\end{definition}
\begin{definition}
	设$m_k$表示类别$k$的某个评价指标，类别$k$的真实样本数为$n_k$，该指标的Weighted平均定义为：
	\begin{equation*}
		m_{\mathrm{weighted}}=\sum_{k=1}^{K}\frac{n_k}{n}m_k
	\end{equation*}
\end{definition}
\begin{definition}
	Micro平均定义为：
	\begin{gather*}
		\operatorname{Precision}_{\mathrm{micro}}=\frac{\sum\limits_{k=1}^{K}\operatorname{TP}_k}{\sum\limits_{k=1}^{K}(\operatorname{TP}_k+\operatorname{FP}_k)} \\
		\operatorname{Recall}_{\mathrm{micro}}=\frac{\sum\limits_{k=1}^{K}\operatorname{TP}_k}{\sum\limits_{k=1}^{K}(\operatorname{TP}_k+\operatorname{FN}_k)} \\
		\operatorname{F1}_{\mathrm{micro}}=2\times\frac{\operatorname{Precision}_{\mathrm{micro}}\operatorname{Recall}_{\mathrm{micro}}}{\operatorname{Precision}_{\mathrm{micro}}+\operatorname{Recall}_{\mathrm{micro}}}
	\end{gather*}
\end{definition}
\begin{note}
	在每个样本只有一个真实标签且模型只预测一个类别的单标签多分类问题中，每个预测正确的样本产生一个真阳性，而每个预测错误的样本同时产生一个假阳性和一个假阴性。因此：
	\begin{equation*}
		\sum_{k=1}^{K}\operatorname{FP}_k=\sum_{k=1}^{K}\operatorname{FN}_k=n-\sum_{k=1}^{K}\operatorname{TP}_k
	\end{equation*}
	从而：
	\begin{equation*}
		\operatorname{Precision}_{\mathrm{micro}}=\operatorname{Recall}_{\mathrm{micro}}=\operatorname{F1}_{\mathrm{micro}}=\operatorname{Accuracy}
	\end{equation*}
	所以在多分类问题中，同时报告Accuracy和Micro通常不会提供额外信息。\par
	若希望所有类别具有相同的重要性，应重点考察Macro平均；若希望评价结果反映实际类别分布，可以使用Weighted平均；若希望从所有样本层面汇总分类结果，则可以使用Micro平均。
\end{note}

\subsection{回归任务}
设真实值为$\{y_i\}_{i=1}^{n}$，模型预测值为$\{\hat{y}_i\}_{i=1}^{n}$。
\begin{definition}
  称：
  \begin{equation*}
    \operatorname{MAE}=\frac{1}{n}\sum_{i=1}^n|y_i-\hat{y}_i|
  \end{equation*}
  为\gls{MAE}。
\end{definition}
\begin{note}
  MAE计算预测误差的绝对值并取平均，因而对所有误差赋予相同的线性权重。它具有易于解释、与原数据同量纲等优点，并对单个异常值不敏感。
\end{note}
\begin{definition}
  称：
  \begin{equation*}
    \operatorname{MAPE}=\frac{1}{n}\sum_{i=1}^n\left|\frac{y_i-\hat{y}_i}{y_i}\right|
  \end{equation*}
  为\gls{MAPE}。
\end{definition}
\begin{note}
  MAPE 将误差按照真实值比例归一化，适合对预测误差的相对大小进行衡量。
\end{note}
\section{模型训练}

\subsection{训练目标与经验风险}
模型训练的目的，是从给定的候选模型中选出一个能够在未知样本上取得良好预测表现的模型。为此，首先需要根据具体任务明确模型表现的评价标准。例如，在分类任务中，可以关注总体分类准确率，也可以更重视某一类别的召回率、精确率或者误分类成本。不同的评价标准体现了不同的任务需求，并决定了哪些预测结果应被认为更好、哪些预测错误应受到更大的惩罚。\par
评价指标主要用于衡量训练完成后的模型是否满足任务需求，但它未必适合直接作为训练时的优化目标。许多评价指标是不连续的、不可微的，或者不能分解为各个样本损失的平均。例如，分类准确率仅取决于最终预测类别是否正确。当模型的预测分数发生小幅变化但预测类别保持不变时，准确率也保持不变，因此它通常无法提供参数应当如何更新的信息。\par
为便于优化，训练时通常使用一个对模型输出变化更加敏感的损失函数作为代理目标。损失函数不必与评价指标具有相同的表达式，但应当与任务目标相容，即损失的减小通常能够促进所关心的预测性能改善。因此，评价指标决定模型最终应达到的目标，损失函数则将这一目标转化为可以进行参数优化的数学问题。\par
设观测数据集$\mathcal{D}=\{(x_i,y_i)\}_{i=1}^{n}$由分布$P$中的$n$个独立同分布样本构成，候选模型组成函数族：
\begin{equation*}
	\mathcal{F}=\{f_\theta:X\to\mathcal{A}:\theta\in\Theta\}
\end{equation*}
其中$\mathcal{A}$为模型输出空间，$\Theta$为参数空间。沿用点估计中的记号，记$L:Y\times\mathcal{A}\to[0,+\infty]$为损失函数，则$L[y,f_\theta(x)]$表示模型$f_\theta$在样本$(x,y)$上产生的预测损失。若分布$P$已知，则模型$f_\theta$在来自同一分布的样本上的平均预测损失可以表示为：
\begin{equation*}
	R(\theta)\coloneq\operatorname{E}(L)
\end{equation*}
称$R(\theta)$为模型$f_\theta$的总体风险。总体风险越小，表示模型在未知样本上产生的平均预测损失越小。因此，从预测角度看，理想的参数应满足：
\begin{equation*}
	\theta^\ast\in\underset{\theta\in\Theta}{\arg\min}R(\theta)
\end{equation*}
然而，分布$P$通常未知，因此总体风险$R(\theta)$无法直接计算。训练时只能利用观测数据，以样本平均损失近似总体风险。
\begin{definition}
	称：
	\begin{equation*}
		\hat{R}(\theta)\coloneq\frac{1}{n}\sum_{i=1}^{n}L[y_i,f_\theta(x_i)]
	\end{equation*}
	为模型$f_\theta$在观测数据集$\mathcal{D}$上的\gls{EmpiricalRisk}。
\end{definition}
以经验风险作为训练目标时，模型训练可以表示为
\begin{equation*}
	\hat{\theta}\in\underset{\theta\in\Theta}{\arg\min}\hat{R}(\theta)
\end{equation*}
这种通过最小化经验风险选择模型参数的方法称为经验风险最小化。由于$\hat{R}(\theta)$依赖于随机抽取的观测数据，因此所得参数$\hat{\theta}$也是数据的函数。
\subsubsection{例子}
下面以分类任务为例，说明如何从模型希望达到的预测效果出发构造损失函数。设类别集合为$Y=\{1,2,\dots,K\}$，分类模型对输入$x$输出一个概率向量：
\begin{equation*}
	q_\theta(x)=(q_{\theta1}(x),q_{\theta22}(),\dots,q_{\theta K}(x))^{\top},\quad
	q_{\theta k}(x)>0,\quad\sum_{k=1}^{K}q_{\theta k}(x)=1
\end{equation*}
其中，$q_{\theta k}(x)$表示模型认为样本$x$属于第$k$类的概率。对于训练样本$(x_i,y_i)$，我们希望模型预测的准确率高，也即赋给真实类别$y_i$的概率$q_{\theta y_i}(x_i)$尽可能接近$1$。\par
由于训练样本相互独立，模型在所有样本上为真实类别给出的概率可以合并为乘积：
\begin{equation*}
	Q_{\mathcal{D}}(\theta)\coloneq\prod_{i=1}^{n}q_{\theta y_i}(x_i)
\end{equation*}
这个乘积越大，说明模型为观测到的真实类别整体赋予的概率越大。最大化$Q_{\mathcal{D}}(\theta)$等价于最小化：
\begin{equation*}
	-\frac{1}{n}\ln Q_{\mathcal{D}}(\theta)=-\frac{1}{n}\sum_{i=1}^{n}\ln q_{\theta y_i}(x_i)
\end{equation*}
由此可以把单个样本的损失定义为：
\begin{equation*}
	L[y_i,q_\theta(x_i)]\coloneq-\ln q_{\theta y_i}(x_i)
\end{equation*}
称该损失为逐样本的\gls{CrossEntropy}损失。当$q_{\theta y_i}(x_i)$接近$1$时，该损失接近$0$；当$q_{\theta y_i}(x_i)$接近$0$时，该损失趋于$+\infty$。因此，模型不仅会因预测类别错误而受到惩罚，还会因确信程度不足而产生较大的损失。相应的经验风险为：
\begin{equation*}
	\hat{R}(\theta)\coloneq\frac{1}{n}\sum_{i=1}^{n}L[y_i,q_\theta(x_i)]=-\frac{1}{n}\sum_{i=1}^{n}\ln q_{\theta y_i}(x_i)
\end{equation*}
%\begin{definition}
%	设 $(X,\mathscr{A},P)$ 为概率空间，$f$为$X$上取值于有限集合
%	$\{x_1,\dots,x_n\}$ 的离散型随机变量，其真实分布为：
%	\begin{equation*}
%		p_i=P(f=x_i),\;i=1,\dots,n
%	\end{equation*}
%	另设$q=(q_1,\dots,q_n)$为定义在同一取值集合上的另一概率分布，满足$q_i>0$。称：
%	\begin{equation*}
%		H(p,q)=-\sum_{i=1}^np_i\ln q_i
%	\end{equation*}
%	为分布$p$相对于分布$q$的\gls{CrossEntropy}。
%\end{definition}
%\begin{property}\label{prop:CrossEntropy}
%	设$p=(p_1,\dots,p_n)$与$q=(q_1,\dots,q_n)$（$p_iq_i\ne0,\;i=1,2,\dots,n$）为定义在同一有限集合上的概率分布，则$H(p,q)\geqslant H(p,p)$，上式取等当且仅当$q_1=p_1,q_2=p_2,\dots,q_n=p_n$ 。
%\end{property}
%\begin{proof}
%	因为：
%	\begin{equation*}
%		H(p,q)=-\sum_{i=1}^np_i\ln q_i=-\sum_{i=1}^np_i\ln\frac{q_i}{p_i}-\sum_{i=1}^np_i\ln p_i
%	\end{equation*}
%	并且由\cref{prop:ConvexFunction}(3)可知$-\ln(x)$在$(0,+\infty)$上为凸函数，由\cref{ineq:Jensen}可知：
%	\begin{equation*}
%		\sum_{i=1}^np_i\ln\frac{q_i}{p_i}\leqslant\ln\left[\sum_{i=1}^{n}p_i\frac{q_i}{p_i}\right]=0
%	\end{equation*}
%	上式取等当且仅当$\dfrac{q_i}{p_i}$为常数，即$q_i=p_i$。所以：
%	\begin{equation*}
%		H(p,q)\geqslant-\sum_{i=1}^np_i\ln p_i=H(p,p)
%	\end{equation*}
%	且当且仅当$q=p$时取等。
%\end{proof}
%交叉熵需要区分真实概率$p$和预测概率$q$。固定真实概率$p=0.25$后，预测概率$q$与$p$一致时交叉熵最小；预测越偏离真实概率，交叉熵越大。
%\begin{figure}[H]
%	\centering
%	\begin{tikzpicture}
%		\begin{axis}[
%			width=9.2cm, height=5.6cm,
%			xlabel={$q$}, ylabel={$H(p,q)$},
%			xmin=0, xmax=1, ymin=0, ymax=4.5,
%			xtick={0,0.25,0.5,0.75,1}, ytick={0,1,2,3,4},
%			axis lines=left, grid=major,
%			grid style={draw=gray!18},
%			samples=200,
%			]
%			\addplot[very thick, orange!85!black, domain=0.02:0.98]
%			{-(0.25*ln(x)+0.75*ln(1-x))};
%			\addplot[only marks, mark=*, orange!85!black] coordinates {(0.25,0.56234)};
%			\draw[dashed, gray] (axis cs:0.25,0) -- (axis cs:0.25,0.56234);
%		\end{axis}
%	\end{tikzpicture}
%	\caption{真实概率固定为$p=0.25$时的二分类交叉熵。横轴$q$为模型预测第一类的概率。}
%\end{figure}
%\begin{note}
%	由\cref{prop:CrossEntropy}可知，交叉熵在所有概率分布$q$中以$q=p$为唯一最小值。实际中真实分布$p$通常未知，故以样本的经验分布$\hat{p}$代替$p$，并最小化$H(\hat{p},q)$。在$H(\hat{p},q)$能够充分逼近$H(p,q)$的条件下，该过程使所估计的分布$q$趋近于真实分布$p$。
%\end{note}
%\begin{note}
%	在分类问题中，我们可以用逐样本的交叉熵作为损失函数。对于第$i$个样本，若其真实类别为第$k$类，则其真实分布$p_i$是一个one-hot向量，即第$k$个分量为$1$，其余分量为$0$。例如，在二分类问题中有：
%	\begin{equation*}
%		p_i=(1,0)^{\top}\qquad\text{或}\qquad p_i=(0,1)^{\top}
%	\end{equation*}
%	若模型对该样本输出概率分布$q_i$，则其交叉熵为：
%	\begin{equation*}
%		H(p_i,q_i)=-\sum_{j=1}^Kp_{ij}\ln q_{ij}=-\ln q_{iy_i}
%	\end{equation*}
%	整个数据集上的交叉熵损失通常取各样本损失的平均：
%	\begin{equation*}
%		L_{\mathrm{CE}}=\frac{1}{n}\sum_{i=1}^nH(p_i,q_i)
%	\end{equation*}
%	由此可见，分类问题中的交叉熵分别比较每个样本的真实分布$p_i$与预测分布$q_i$。由于$p_i$是one-hot向量，$H(p_i,q_i)$实际上只取决于模型赋给该样本真实类别的概率$q_{iy_i}$：该概率越接近$1$，损失越小；该概率越接近$0$，损失越大。
%\end{note}

\subsection{过拟合}
经验风险$\hat{R}(\theta)$只度量模型在观测数据集上的平均损失，总体风险$R(\theta)$则度量模型对同一分布中新样本的平均损失，二者通常并不相等。\par
有时持续降低$\hat{R}(\theta)$未必会同时降低$R(\theta)$：模型可能不是在提取能够推广到新样本的稳定规律，而是利用观测数据中的偶然噪声、异常点甚至样本的细微特征，对有限训练集作近似记忆。此时训练损失较低，而新样本上的损失较高，这种现象称为\gls{Overfitting}。\par
\begin{figure}[htbp]
	\centering
	\makebox[\textwidth][c]{%
		\begin{tikzpicture}
			\begin{groupplot}[
				group style={
					group size=3 by 1,
					horizontal sep=0.025\textwidth
				},
				width=0.30\textwidth,
				height=0.27\textwidth,
				xmin=-1.25,
				xmax=1.25,
				ymin=-2.2,
				ymax=2.2,
				domain=-1.2:1.2,
				samples=300,
				axis lines=middle,
				xlabel={$x$},
				ylabel={$y$},
				xtick={-1,0,1},
				ytick={-2,-1,0,1,2},
				tick label style={font=\small},
				label style={font=\small},
				title style={
					font=\small,
					yshift=-0.5ex
				},
				clip=false
				]
				
				% 第一种复杂拟合
				\nextgroupplot[
				title={Cubic fit}
				]
				
				\addplot[
				dashed,
				thick,
				gray
				] {x};
				
				\addplot[
				thick
				] {
					x
					+1.8*(x+1)*x*(x-1)
				};
				
				\addplot[
				only marks,
				mark=*,
				mark size=2.2pt
				] coordinates {
					(-1,-1)
					(0,0)
					(1,1)
				};
				
				% 第二种复杂拟合
				\nextgroupplot[
				title={Quintic fit},
				legend style={
					at={(0.5,-0.34)},
					anchor=north,
					legend columns=-1,
					draw=none,
					font=\small,
					/tikz/every even column/.append style={column sep=0.8em}
				}
				]
				
				\addplot[
				dashed,
				thick,
				gray
				] {x};
				\addlegendentry{true relation $y=x$}
				
				\addplot[
				thick
				] {
					x
					-3.2*(x+1)*x*(x-1)*(x^2-0.35)
				};
				\addlegendentry{fitted curve}
				
				\addplot[
				only marks,
				mark=*,
				mark size=2.2pt
				] coordinates {
					(-1,-1)
					(0,0)
					(1,1)
				};
				\addlegendentry{training point}
				
				% 第三种复杂拟合
				\nextgroupplot[
				title={Oscillatory fit}
				]
				
				\addplot[
				dashed,
				thick,
				gray
				] {x};
				
				\addplot[
				thick
				] {
					x
					+1.1*(x+1)*x*(x-1)
					*sin(deg(7*x))
				};
				
				\addplot[
				only marks,
				mark=*,
				mark size=2.2pt
				] coordinates {
					(-1,-1)
					(0,0)
					(1,1)
				};
				
			\end{groupplot}
			
		\end{tikzpicture}%
	}
	\caption{取自$y=x$的三个样本可由不同的模型进行解释。}
	\label{fig:overfitting-three-points}
\end{figure}
从\cref{fig:overfitting-three-points}可以看出，三条复杂程度不同的曲线都完美地拟合了由三个观测点构成的观测数据集，因此它们在该观测集上取得了相同且最小的经验风险$\hat{R}(\theta)$。然而，在训练样本之外，若新的观测仍来自真实关系$y=x$，这些模型会产生较大的预测误差，即它们存在过拟合的现象。由此可见，仅凭观测数据上的经验风险无法判断模型的真实预测能力；即使多个模型具有相同的训练误差，它们的总体风险$R(\theta)$也可能存在显著差异。\par
综上，我们需要关注差值：
\begin{equation*}
	R(\theta)-\hat{R}(\theta)
\end{equation*}
它反映了模型从观测数据推广到新数据时的性能差异。
\subsubsection{训练集、验证集与测试集}
计算上述差值需要计算$R(\theta)$，但是由于分布$P$未知，$R(\theta)$通常无法直接计算，只能通过未参与训练的数据加以估计，所以我们通常将原始数据随机划分为三个互不相交的部分：
\begin{equation*}
	\mathcal{D}=\mathcal{D}_{\mathrm{train}}\cup\mathcal{D}_{\mathrm{val}}\cup\mathcal{D}_{\mathrm{test}}
\end{equation*}
分别称为\gls{TrainingSet}、\gls{ValidationSet}和\gls{TestSet}。三者应当尽可能来自相同的目标分布；对于类别不平衡的分类问题，通常采用分层抽样，使各部分的类别比例大致一致。\par
\begin{method}
	数据集划分后的模型训练与选择通常按下列顺序进行：
	\begin{enumerate}
		\item 在训练集$\mathcal{D}_{\mathrm{train}}$上估计模型的可学习参数$\theta$；
		\item 在验证集$\mathcal{D}_{\mathrm{val}}$上比较不同模型的性能，并据此选择模型；
		\item 确定全部建模方案后，只在测试集$\mathcal{D}_{\mathrm{test}}$上进行一次最终评价，报告模型对未见数据的性能估计。
	\end{enumerate}
\end{method}
验证集虽然不参与参数的直接拟合，却参与模型选择；若反复根据测试集结果修改模型，测试集实际上也参与了训练流程，其评价结果便会过于乐观。因此，测试集必须与整个训练和选择过程隔离。数据预处理也应遵循相同原则：缺失值填补、标准化、特征选择等步骤只能在训练集上估计所需参数，再将已经确定的变换应用到验证集和测试集，不能提前利用后两者的信息。\par
当样本量较小时，如果再从训练数据中单独划分一个固定的验证集会进一步减少模型实际用于训练的样本。此时可以在训练数据内部使用$K$折\gls{CV}：将训练数据近似均匀地分为$K$份，依次将其中一份作为验证集，其余$K-1$份作为训练集，共进行$K$轮训练与验证，并对$K$次验证结果取平均，以评价候选模型。此时仍应保留独立测试集；交叉验证用于模型选择，测试集只用于最终评价。
\subsubsection{正则化与模型空间约束}
对$f=f_\theta\in\mathcal{F}$，记$\hat{R}_{\mathrm{train}}(f)\coloneq\hat{R}_{\mathrm{train}}(\theta)$及$R(f)\coloneq R(\theta)$。\par
若$\mathcal{F}_1\subseteq\mathcal{F}_2$，则：
\begin{equation*}
	\inf_{f\in\mathcal{F}_2}\hat{R}_{\mathrm{train}}(f)\leqslant\inf_{f\in\mathcal{F}_1}\hat{R}_{\mathrm{train}}(f)
\end{equation*}
也就是说，扩大模型空间不会使最优训练风险变大，因为新增的模型为拟合训练数据提供了更多选择；但总体风险并不具有相同的单调性。模型空间越大，训练过程越容易从大量候选模型中找到一个恰好迎合当前样本偶然性的模型，因而经验风险虽然可以很低，经验风险与总体风险之间的偏差却更难控制。极端情况下，一个足够复杂的模型可以记住每个训练样本，使训练风险接近$0$；但训练样本之外的函数取值仍可能几乎不受约束，这正是过拟合产生的原因。
\begin{definition}
	设$\Omega:\mathcal{F}\to[0,+\infty]$为模型复杂度函数，$\Omega(f)$越大表示模型$f$越复杂。对$c\geqslant0$，定义受限模型空间：
	\begin{equation*}
		\mathcal{F}_c\coloneq\{f\in\mathcal{F}:\Omega(f)\leqslant c\}
	\end{equation*}
	在$\mathcal{F}_c$中最小化经验风险：
	\begin{equation*}
		\hat{f}_c\in\underset{f\in\mathcal{F}_c}{\arg\min}\hat{R}_{\mathrm{train}}(f)
	\end{equation*}
	称为约束形式的\gls{Regularization}。其中，$c$称为复杂度上界。
\end{definition}
实际计算中更常使用正则化的惩罚形式：
\begin{equation*}
	\hat{f}_\lambda\in\underset{f\in\mathcal{F}}{\arg\min}\left\{\hat{R}_{\mathrm{train}}(f)+\lambda\Omega(f)\right\},\;\lambda\geqslant0
\end{equation*}
其中，$\hat{R}_{\mathrm{train}}(f)+\lambda\Omega(f)$称为正则化目标函数，$\lambda$称为\gls{RegularizationStrength}。当$\lambda=0$时，该问题退化为原来的经验风险最小化；$\lambda$越大，模型复杂度在目标函数中付出的代价越大，训练过程便越倾向于选择简单模型。
\begin{note}
	限制模型空间的本质是在拟合训练数据与控制模型复杂度避免过拟合之间取得平衡：模型空间过小会导致模型无法很好的拟合训练数据，也即\gls{Underfitting}，模型空间过大会导致过拟合。
\end{note}

\input{statistics/machine-learning/tree}
\input{statistics/machine-learning/knn}
\section{核方法}

\begin{definition}
	设$X$是空间，$k:X\times X\to\mathbb{R}^{}$被称为\gls{KernelFunction}。
	\begin{enumerate}
		\item 若对任意的$x_1,x_2\in X$有$k(x_1,x_2)=k(x_2,x_1)$，则称$k$是\textbf{对称核}；
		\item 若对任意的$\seq{x}{n}\in X$和任意的$\seq{a}{n}\in\mathbb{R}^{}$有：
		\begin{equation*}
			\sum_{i=1}^{n}\sum_{j=1}^{n}a_ia_jk(x_i,x_j)\geqslant0
		\end{equation*}
		则称$k$为\textbf{半正定核}。
	\end{enumerate}
\end{definition}
\begin{definition}
	设$X$是空间，$k:X\times X\to\mathbb{R}^{}$。对任意的$\seq{x}{n}\in X$，称$K=\Big(k(x_i,x_j)\Big)\in M_{n\times n}(\mathbb{R}^{})$为$n$维\textbf{核矩阵}或\textbf{Gram矩阵}。
\end{definition}
\begin{definition}
	设$X$是空间，$k:X\times X\to\mathbb{R}^{}$，$Y$是Hilbert空间。若存在$\varphi:X\to Y$使得对任意的$x_1,x_2\in X$有：
	\begin{equation*}
		k(x_1,x_2)=\Big(\varphi(x_1),\varphi(x_2)\Big)
	\end{equation*}
	则称$\varphi$是$k$关于$Y$的一个\textbf{特征映射}。
\end{definition}
\begin{theorem}\label{theo:FeatureMapPositiveSemidefiniteKernel}
	设$X$是空间，$k:X\times X\to\mathbb{R}^{}$，$Y$是Hilbert空间。若$k$存在关于$Y$的特征映射，则$k$是半正定核。
\end{theorem}
\begin{proof}
	对任意的$\seq{x}{n}\in X$和任意的$\seq{a}{n}\in\mathbb{R}^{}$有：
	\begin{align*}
		\sum_{i=1}^{n}\sum_{j=1}^{n}a_ia_jk(x_i,x_j)&=\sum_{i=1}^{n}\sum_{j=1}^{n}a_ia_j\Big(\varphi(x_i),\varphi(x_j)\Big)=\left(\sum_{i=1}^{n}a_i\varphi(x_i),\sum_{i=1}^{n}a_j\varphi(x_i)\right) \\
		&=\left\|\sum_{i=1}^{n}a_i\varphi(x_i)\right\|^2\geqslant0\qedhere
	\end{align*}
\end{proof}
\begin{definition}
	若Hilbert空间$Y$是定义在$X$上的实值函数空间，则称$Y$为$X$上的\textbf{Hilbert函数空间}。
\end{definition}
\begin{definition}
	设$Y$是$X$上的实Hilbert函数空间。若对任意的$x\in X$，泛函：
	\begin{equation*}
		L_x:Y\to\mathbb{R}^{},\quad\forall\;f\in Y,\;L_x(f)=f(x)
	\end{equation*}
	都是有界线性泛函，则称$Y$为$X$上的\gls{RKHS}。
\end{definition}
\begin{theorem}\label{theo:ReproducingKernel}
	若$Y$是$X$上的RKHS，则对于任意的$x\in X$，存在唯一的$k_x\in Y$使得对于任意的$f\in Y$有$f(x)=(f,k_x)$。
\end{theorem}
\begin{proof}
	因为$Y$是$X$上的RKHS，所以对于任意的$x\in X$，泛函：
	\begin{equation*}
		L_x:Y\to\mathbb{R}^{},\quad\forall\;f\in Y,\;L_x(f)=f(x)
	\end{equation*}
	都是有界线性泛函。由\cref{theo:LinearOperatorBoundedContinuous}(1)可知$L_x$是连续泛函，根据\cref{theo:RieszRepresentation}可知存在唯一的$k_x\in Y$使得对任意的$f\in Y$有$L_x(f)=f(x)=(f,k_x)$。
\end{proof}
\begin{definition}
	设$Y$是$X$上的RKHS，$x\in X$。称满足对于任意的$f\in Y$有$f(x)=(f,k_x)$的$k_x\in Y$为$x$关于$Y$的\gls{ReproducingKernel}。
\end{definition}
\begin{note}
	再生核族$\{k_x:x\in X\}$可导出核函数$k(x_1,x_2)=k_{x_1}(x_2)$。
\end{note}
\begin{property}\label{prop:ReproducingKernel}
	设$Y$是$X$上的RKHS，$k$是$Y$上的再生核族$\{k_x:x\in X\}$导出的核函数，则：
	\begin{enumerate}
		\item 对任意的$x_1,x_2\in X$有$k(x_1,x_2)=(k_{x_1},k_{x_2})$；
		\item $k$是对称核；
		\item $k$是半正定核；
	\end{enumerate}
\end{property}
\begin{proof}
	(1)取任意的$x_1,x_2\in X$，由\cref{theo:ReproducingKernel}可知$k_{x_1},k_{x_2}\in Y$，所以$k_{x_1}(x_2)=k(x_1,x_2)=(k_{x_1},k_{x_2})$。\par
	(2)在(1)中其实有$k_{x_1}(x_2)=k(x_1,x_2)=(k_{x_1},k_{x_2})=(k_{x_2},k_{x_1})=k_{x_2}(x_1)=k(x_2,x_1)$。\par
	(3)由\cref{theo:FeatureMapPositiveSemidefiniteKernel}立即可得。
\end{proof}
\begin{theorem}[Aronszajn Theorem]\label{theo:AronszajnTheorem}
	设$X$是空间，$k:X\times X\to\mathbb{R}^{}$。$k$是半正定核当且仅当存在$X$上的再生核Hilbert空间$Y$，对任意的$x\in X$，$k(x,\cdot)$是$x$关于$Y$的再生核。
\end{theorem}
\begin{proof}
	\textbf{(1)充分性：}由\cref{prop:ReproducingKernel}(2)立即可得。
%	\textbf{(2)必要性：}记：
%	\begin{equation*}
%		Y_0=\left\{f(x)=\sum_{i=1}^{n}a_ik(x,x_i):\forall\;n\in\mathbb{N}^+,\;a_i\in\mathbb{R}^{},\;x_i\in X\right\}
%	\end{equation*}
%	对任意的：
%	\begin{equation*}
%		f(x)=\sum_{i=1}^{m}a_ik(x,x_i),g(x)=\sum_{i=1}^{n}b_ik(x,y_i)\in Y_0
%	\end{equation*}
%	定义：
%	\begin{equation*}
%		(f,g)=\sum_{i=1}^{m}\sum_{j=1}^{n}a_ib_jk(x_i,y_j)
%	\end{equation*}
%	则$(\cdot,\cdot)$是内积（线性性和对称性显然，由$k$的半正定性可得非负性），于是$Y_0$成为一个内积空间。由\info{度量空间的完备化}可将$Y_0$完备化得到Hilbert空间$Y$
\end{proof}
\begin{theorem}[Representor Theorem]\label{theo:RepresentorTheorem}
	设$X$是空间，$k:X\times X\to\mathbb{R}^{}$是半正定核，由\cref{theo:AronszajnTheorem}，记$Y$是由$k$导出的$X$上的RKHS，$\{(x_i,y_i)\}_{i=1}^{n}$为训练样本。对于优化问题：
	\begin{equation*}
		\min_{f\in Y}\Omega[f(x_1),f(x_2),\dots,f(x_n)]+\lambda||f||^2,\quad\lambda>0
	\end{equation*}
	其任一最优解$f^*$都可以表示为：
	\begin{equation*}
		f(x)=\sum_{i=1}^{n}a_ik(x,x_i)
	\end{equation*}
	其中$a_i\in\mathbb{R}^{}$。
\end{theorem}
\begin{proof}
	由\cref{prop:SpanSubspace}(1)，取$Y$的子空间：
	\begin{equation*}
		Y_n=\operatorname{span}\{k_{x_1},k_{x_2},\dots,k_{x_n}\}
	\end{equation*}
	由\cref{prop:SpanSubspace}(2)、\cref{prop:FiniteDimensionalNormedLinearSpace}(2)和\cref{prop:CauchySeq}(1)可知$Y_n$是闭的。根据\cref{theo:OrthogonalProjection}可知对任意的$f\in Y$，存在下述唯一的分解：
	\begin{equation*}
		f=f_1+f_2,\quad f_1\in Y_n,\;f_2\in Y_n^{\perp}
	\end{equation*}
	由\cref{prop:Orthogonal}(1)可知$||f||^2=||f_1||^2+||f_2||^2$。因为$f_2\in Y_n^{\perp}$，所以对任意的$i=1,2,\dots,n$有$f_2\perp k_{x_i}$，于是：
	\begin{equation*}
		(f_2,k_{x_i})=f_2(x_i)=0
	\end{equation*}
	所以$f(x_i)=f_1(x_i)$，优化问题可被改写为：
	\begin{equation*}
		\min_{f\in Y}\Omega[f_1(x_1),f_1(x_2),\dots,f_1(x_n)]+\lambda[||f_1||^2+||f_2||^2],\quad\lambda>0
	\end{equation*}
	若$f^*$是最优解，则$f^*$必然在$Y_n$中，否则$f^*$在$Y_n$上的正交投影对应的目标函数值一定比$f^*$对应的目标函数值小。
\end{proof}
\input{statistics/machine-learning/nn}
\section{不平衡样本集问题}

\begin{definition}
	在分类问题中，若观测数据集中不同类别的样本数量相差悬殊，则称该数据集存在\gls{Class-imbalancedDataset}问题。此时，样本数量较多的类别称为\gls{MajorityClass}，样本数量较少的类别称为\gls{MinorityClass}。
\end{definition}
不平衡样本集问题的根源在于，模型训练通常以最小化经验风险为目标，而经验风险是各观测样本损失的平均。当不同类别的样本数量相差悬殊时，多数类样本对经验风险的贡献会占据主导地位，因此，即使模型在少数类上的预测表现很差，只要能够较好地拟合多数类，仍可能获得较小的总体经验风险。换言之，经验风险最小化本身并不会保证模型对各类别具有相近的识别能力，而可能使训练过程系统性地偏向多数类。\par
例如，在一个二分类任务中，观测数据集包含$100$个样本，其中$99$个属于多数类$0$，仅有$1$个属于少数类$1$。若模型对任意样本的输出都为$0$，则在$0$-$1$损失下（预测正确损失函数值为$0$，预测错误损失函数值为$1$），其经验风险仅为$0.01$。因此，从总体经验风险的角度看，该模型已经取得了看似良好的预测结果。然而，它完全无法识别任何少数类样本。\par
类别不平衡现象广泛存在于实际分类任务中，并且在许多应用场景下，样本数量较少的类别恰恰对应我们真正关心的事件。例如，疾病患者、欺诈交易和设备故障通常只占全部观测中的很小比例，但识别这些少数类样本往往正是建模任务的核心目标。因此，一个仅通过良好预测多数类而获得较低经验风险的模型，即使在整体意义上表现良好，也可能完全失去实际应用价值，因为它没有学习到任务中最重要的少数类模式。\par
在接下来的讨论中，我们将主要涉及二分类任务，并且默认将少数类记为正类。

\subsection{评价指标}
类别不平衡不仅会影响模型的训练过程，也会影响模型性能的评价。当不同类别的样本数量相差悬殊时，一些在一般分类任务中常用的评价指标可能无法准确反映模型对少数类的识别能力。因此，在讨论不平衡样本集的评价问题时，一个自然的思路是重新考察此前介绍的各类评价指标，分析它们在类别不平衡情形下是否仍然具有合理的解释。对于容易受到类别比例影响的指标，则需要考虑相应的修正形式或采用新的评价指标。\par
在类别不平衡的数据集中，准确率显然容易受到多数类样本的主导。为避免不同类别的样本数量直接决定评价结果，可以分别计算模型在各类别上的识别准确率，再对其进行等权平均，从而减弱类别数量不平衡对评价结果的影响。。
\begin{definition}
	对于二分类任务，称：
	\begin{equation*}
		\operatorname{BalancedAccuracy}=\frac{1}{2}\left(\frac{\operatorname{TP}}{\operatorname{TP}+\operatorname{FN}}+\frac{\operatorname{TN}}{\operatorname{TN}+\operatorname{FP}}\right)
	\end{equation*}
	为模型的\gls{BalancedAccuracy}。
\end{definition}
由于默认将少数类视为正类，因此精确率、召回率、$\operatorname{F1}$分数、$\operatorname{F}_{\beta}$分数以及$\operatorname{AUPRC}$都能够直接反映模型在少数类上的预测表现。精确率刻画模型预测为少数类的样本中有多少真正属于少数类，召回率刻画实际少数类样本中有多少被模型成功识别，$\operatorname{F1}$分数和$\operatorname{F}_{\beta}$分数则综合衡量精确率与召回率，而$\operatorname{AUPRC}$进一步描述不同分类阈值下精确率与召回率之间的整体权衡。这些指标更加关注正类的识别效果，因此通常适用于类别不平衡问题。

\subsection{代价敏感学习}
对于不平衡分类问题，一种直接的处理思路是重新调整不同类别样本在损失函数中的权重，以抵消类别样本数量差异对经验风险的影响，使各类别的预测误差能够在训练目标中得到更为均衡的体现。这类通过赋予不同类型样本的预测错误以不同代价，从而改变模型优化目标的方法称为\gls{CostSensitiveLearning}。\par
具体而言，设分类问题共有$K$个类别，第$k$类样本数量为$n_k$，为该类别赋予权重$w_k>0$，则代价敏感学习一般的经验风险被定义为：
\begin{equation*}
	\hat{R}_w(\theta)\coloneq\sum_{i=1}^{n}w_{y_i}L[y_i,f_\theta(x_i)]
\end{equation*}\par
一种常用的权重设定方式是令类别权重与该类别的样本数量成反比，即：
\begin{equation*}
	w_k\propto\frac{1}{n_k}
\end{equation*}
此时第$k$类样本在加权经验风险中的总权重满足：
\begin{equation*}
	n_kw_k\propto1
\end{equation*}
因此，不同类别虽然具有不同的样本数量和单样本权重，但其在总体损失中具有相同的总权重，从而消除了类别频率本身对经验风险贡献所造成的失衡。

\subsection{样本平衡类方法}
与代价敏感学习直接调整不同类别样本在损失函数中的贡献不同，处理不平衡样本集的另一类基本思路是从训练数据本身出发，通过改变不同类别样本在训练过程中的出现频率，使模型所观察到的经验类别分布更加均衡。这类方法统称为\textbf{样本平衡类方法}。其核心思想是，在不改变原始学习任务和损失函数形式的前提下，通过对多数类样本进行删减、对少数类样本进行重复或合成，重新构造用于模型训练的数据集，从而减弱多数类样本因数量占优而对经验风险及参数更新产生的主导作用。
\subsubsection{欠采样与过采样}
设分类问题共有$K$个类别，第$k$类样本数量为$n_k$。样本平衡最直接的做法是通过重新采样改变不同类别样本参与模型训练的频率，从而调整训练数据中的经验类别分布。根据调整方式的不同，可以分为\gls{UnderSampling}与\gls{OverSampling}。\par
欠采样通过减少多数类样本的数量来缓解类别之间的样本规模差异。最简单的随机欠采样方法从多数类样本中随机选取部分观测用于训练。对于二分类问题，设多数类和少数类的样本数量分别为$n_{\mathrm{maj}}$和$n_{\mathrm{min}}$，可以从多数类中抽取$n_{\mathrm{maj}}'<n_{\mathrm{maj}}$个样本，使$n_{\mathrm{maj}}'$与$n_{\mathrm{min}}$之间的差异缩小。\par
欠采样能够减少多数类对训练目标的主导作用，同时直接降低训练数据规模，因此在多数类样本数量极大时还具有降低存储与计算成本的优势。其主要缺点在于部分已经观测到的样本会被舍弃。如果多数类具有复杂的内部结构，或者被删除的样本包含重要的决策边界信息，则过强的欠采样可能造成有效信息损失，并降低模型的泛化能力。\par
与之相对，过采样不删除原有观测，而是通过提高少数类样本在训练过程中的出现频率来改善类别不平衡。最简单的随机过采样方法从少数类样本中进行有放回抽样，使其有效样本数量增加。例如，可以将少数类从$n_{\mathrm{min}}$个样本重复采样至$n_{\mathrm{min}}'$个，使其与多数类的样本数量更加接近。\par
随机过采样避免了欠采样所带来的信息丢失，但简单复制已有样本并不会产生新的观测信息。少数类中的同一观测可能被模型反复使用，从而增加对有限少数类样本过拟合的风险；与此同时，重复采样还会增大训练数据规模和计算开销。
\subsubsection{SMOTE}
随机过采样通过重复已有少数类样本提高其在训练过程中的出现频率，但重复观测本身并不会引入新的特征信息。为缓解这一问题，\gls{SMOTE}不再直接复制少数类样本，而是利用少数类样本在特征空间中的局部邻域结构生成新的合成样本。\par
设$x_i$为某个少数类样本，并在少数类样本中寻找其$k$个近邻。从这些近邻中选取一个样本$x_j$，再生成随机变量：
\begin{equation*}
	\lambda\sim\operatorname{Uniform}(0,1)
\end{equation*}
则SMOTE生成的新样本为：
\begin{equation*}
	x_{\mathrm{new}}=x_i+\lambda(x_j-x_i)
\end{equation*}
因此，$x_{\mathrm{new}}$位于$x_i$与$x_j$之间的线段上，并被赋予与$x_i$相同的类别标签。通过对不同的少数类样本及其近邻重复这一过程，可以逐步扩充少数类训练样本。\par
SMOTE隐含的基本假设是：相邻少数类样本之间的区域仍然属于少数类，因此可以通过局部线性插值生成新的少数类样本。相较于随机过采样直接复制已有观测，这种方式能够扩展少数类样本在特征空间中的覆盖范围。这一假设也是SMOTE的主要局限。当少数类靠近决策边界或少数类本身具有复杂的非线性结构时，两个少数类样本之间的插值区域未必仍然属于少数类，从而可能生成不合理的合成样本并干扰决策边界。

\subsection{Balanced Bagging}
Bagging通过对训练数据进行重复采样并训练多个基学习器，再将各基学习器的预测结果进行集成，利用基学习器之间的异质性去提高模型的预测精度（\cref{prop:TreeEnsemble}(2)）。对于不平衡分类问题，可以进一步将类别平衡操作嵌入Bagging的采样过程中，使每个基学习器均在相对平衡的数据集上训练。这类方法通常称为\textbf{Balanced Bagging}。\par
设二分类训练集为：
\begin{equation*}
	\mathcal{D}=\mathcal{D}_{+}\cup\mathcal{D}_{-}
\end{equation*}
其中$n_{+}\ll n_{-}$。对于第$b$个基学习器，Balanced Bagging保留少数类样本，并从多数类中随机抽取与少数类数量相近的样本，构造平衡训练集$\mathcal{D}^{(b)}$并训练$f^{(b)}$，其最终模型为$f^{(1)},f^{(2)},\dots,f^{(B)}$经过多数投票得到的结果。\par
这种方法通常采用多数类欠采样，而不是少数类过采样。原因在于单次欠采样虽然会丢失多数类信息，但Bagging中的不同基学习器可以使用不同的多数类子集，从而在整个集成模型中逐步利用更多多数类观测，因此Bagging恰好能够缓解欠采样最主要的信息损失问题。相比之下，随机过采样主要是在不同训练集中重复有限的少数类样本，并不会提供新的观测信息，还会增大每个基学习器的训练规模，并可能使不同基学习器反复依赖相同的少数类样本，提高了基学习器之间的相关性，从而无法达到理想的预测精度。

\subsection{困难负样本挖掘}
\begin{definition}
	设二分类训练集为$\mathcal{D}=\mathcal{D}_{+}\cup\mathcal{D}_{-}$，其中$\mathcal{D}_{+}$与$\mathcal{D}_{-}$分别为正类和负类样本集。对于给定模型$f_{\theta}$，定义负类样本$(x_i,0)\in\mathcal{D}_{-}$的分类困难程度为：
	\begin{equation*}
		d_i(\theta)\coloneq L[0,f_{\theta}(x_i)]
	\end{equation*}
	给定阈值$\tau$，若$d_i(\theta)>\tau$，则称$(x_i,0)$为模型$f_{\theta}$的困难负样本。从多数类中选择困难负样本并用于后续训练的方法称为\gls{HardNegativeMining}。
\end{definition}
困难负样本挖掘的核心思想是：多数类样本虽然数量庞大，但并非所有样本对模型训练都具有同等价值。大量能够被当前模型轻易正确分类的负类样本对决策边界的进一步修正作用有限，而那些容易被误判、损失较大的负类样本更能反映模型当前尚未解决的分类困难。因此，与随机欠采样等概率地丢弃多数类样本不同，困难负样本挖掘利用当前模型主动筛选更有训练价值的负类样本，使模型将有限的训练能力集中于难以区分的区域。\par
由于样本的困难程度由当前模型决定，困难负样本并不是固定不变的。随着模型参数更新，原有困难负样本的损失可能降低，而其他负类样本可能成为新的困难样本。因此，困难负样本挖掘通常与模型训练交替进行，即利用当前模型筛选困难负样本，再使用这些样本更新模型，并重复这一过程，直至满足停止准则。
\subsubsection{机器学习中的困难负样本挖掘}
\begin{method}
	设第$t$轮训练所使用的负类样本集为$\mathcal{S}_{-}^{(t)}\subseteq\mathcal{D}_{-}$。对于非神经网络的传统机器学习模型，困难负样本挖掘的一般过程为：
	\begin{enumerate}
		\item 从$\mathcal{D}_{-}$中选取部分负类样本，构造初始负类训练集$\mathcal{S}_{-}^{(0)}$；
		\item 对$t=0,1,2,\ldots$，使用$\mathcal{D}_{+}\cup\mathcal{S}_{-}^{(t)}$训练模型$f_{\theta^{(t)}}$；
		\item 使用$f_{\theta^{(t)}}$评估负类样本的困难程度$d_i(\theta^{(t)})$，并据此确定当前模型下需要重点用于后续训练的困难负样本；
		\item 根据当前负类训练集以及新识别出的困难负样本构造下一轮负类训练集$\mathcal{S}_{-}^{(t+1)}$，随后进入下一轮训练。
	\end{enumerate}
	当验证集上的评价指标不再明显改善，或达到预设的最大迭代轮数时停止。
\end{method}
设$n_{+}=|\mathcal{D}_{+}|$且$s_t=|\mathcal{S}_{-}^{(t)}|$。若还希望正类与选中负类对训练目标具有相同的总贡献，可以在第$t$轮使用如下经验风险：
\begin{equation*}
	\hat{R}_{\mathrm{HNM}}^{(t)}(\theta)\coloneq\frac{1}{n_{+}}\sum_{(x_i,1)\in\mathcal{D}_{+}}L[1,f_{\theta}(x_i)]
	+\frac{1}{s_t}\sum_{(x_i,0)\in\mathcal{S}_{-}^{(t)}}L[0,f_{\theta}(x_i)]
\end{equation*}\par
负类训练集$\mathcal{S}_{-}^{(t)}$的更新方式并不唯一。最直接的做法是在已有负类训练集上不断加入新识别出的困难负样本，使训练集随迭代逐渐扩增；也可以在每一轮根据当前模型重新选择负类样本，从而完全重构$\mathcal{S}_{-}^{(t+1)}$。此外，还可以采用介于两者之间的策略，例如保留部分已有困难负样本，同时移除已经能够被模型轻易识别的样本并补充新的困难负样本。具体采用何种更新方式取决于具体场景。
\subsubsection{深度学习中的困难负样本挖掘}
设模型给出正类概率$q_{\theta}(x_i)\in(0,1)$，记：
\begin{equation*}
	p_{ti}=
	\begin{cases}
		q_{\theta}(x_i), &y_i=1 \\
		1-q_{\theta}(x_i), &y_i=0
	\end{cases}
	\quad
	\alpha_{ti}=
	\begin{cases}
		\alpha,& y_i=1, \\
		1-\alpha,& y_i=0
	\end{cases}
\end{equation*}
其中$\alpha\in(0,1)$为类别权重参数。
\begin{definition}
	设聚焦参数$\gamma\geqslant0$。称：
	\begin{equation*}
		L_{\mathrm{FL}}[y_i,q_{\theta}(x_i)]
		\coloneq-\alpha_{ti}(1-p_{ti})^{\gamma}\ln p_{ti}
	\end{equation*}
	为样本$(x_i,y_i)$的\gls{FocalLoss}。
\end{definition}
Focal loss在交叉熵前加入调节因子$(1-p_{ti})^{\gamma}$。当样本容易分类时$p_{ti}$接近$1$，其损失被显著抑制；当样本难以分类时$p_{ti}$较小，其损失得到保留。因此，Focal loss无需显式筛除容易样本，即可使训练更加关注困难样本。
\subsubsection{困难负样本挖掘的缺陷}
困难负样本挖掘依赖当前模型对样本困难程度的判断，因此其效果容易受到初始模型和迭代过程的影响。如果模型早期的决策边界存在较大偏差，所选出的困难负样本也可能缺乏代表性，并进一步影响后续训练。此外，高损失样本并不一定是真正有信息的困难样本。错误标签、异常观测以及类别重叠区域中的样本同样可能产生较大损失，从而被反复选中并获得过高的训练关注。与此同时，过度强调困难负样本会削弱模型对多数类整体分布的利用，若负样本选择过于集中，还可能降低模型对其他多数类区域的泛化能力。


\input{statistics/machine-learning/hyperparam}